\subsection{National account trajectories and source-specific forest extent}\label{sec:nationalresults}
Of the 183 eligible account identifiers, total biocapacity increases in 91 and decreases in 92 during 2015--2023. Among increasing totals, the forest account component decreases in 68, increases in 22, and is unresolved at the rounding scale in one. This is account composition, not 68 independently observed forest losses. Domain choice matters: Canada and Cameroon change total-account direction under some component exclusions. All eight candidates, not only nominally discordant ones, remain in the comparison (Table~\ref{tab:nationalcases}).

\begin{table}[htbp]
\centering\small
\caption{2015--2023 national account changes and additional physical evidence. Percentages in the first two columns refer to gha, not physical forest stocks. The last column identifies the evidence type; none supplies the original calibrated national endpoint AGB comparison. Country groups and inventory windows differ.}\label{tab:nationalcases}
\begin{tabularx}{\textwidth}{lrrX}
\toprule
Country & Total (\%) & Forest account (\%) & Additional evidence\\\midrule
Brazil & $-1.21$ & $-3.15$ & Mapped formations and annual transitions\\
Indonesia & $+5.46$ & $-5.92$ & Mapped forest and annual transitions\\
DR Congo & $-2.91$ & $-7.64$ & Source-specific mapped forest series\\
Cameroon & $+0.54$ & $-3.16$ & No completed physical numerical module\\
Canada & $+0.08$ & $-1.10$ & Earlier inventory-domain biomass/change tables\\
Finland & $-1.60$ & $-1.01$ & Inventory-cycle biomass and volume\\
France & $-4.47$ & $+2.91$ & Production-forest inventory-cycle volume\\
Australia & $+12.95$ & $-0.33$ & No completed physical numerical module\\
\bottomrule
\end{tabularx}
\end{table}

Brazilian formations decline from 412.016 to 401.154 million mapped ha, a decrease of 10.863 million ha or 2.64\%. The separate formations-plus-savanna group decreases by 18.184 million ha or 3.43\%. Indonesia's source forest decreases from 107.055 to 106.352 million ha, a net decrease of 702,562 ha or 0.66\%, while mapped oil palm increases by 2.833 million ha. Thus Indonesia has a nominal total-account increase and mapped-forest decrease, whereas Brazil has both decreasing. These are not paired calibrated AGB signs, and their two-case pattern is not a reporting-error rate. The Indonesian earlier-collection forest decline is 3.106 million ha, but its footprint is 43,130 ha larger; the collection-and-support sensitivity is substantial.

Congo's official country table assigns 163.948 and 160.645 million ha to its forest group in 2015 and 2023, a decline of 3.303 million ha or 2.01\%. Dense forest decreases by 1.754 million ha and woodland by 1.577 million ha; mapped mangrove increases by 28,328 ha. Agriculture/pasture increases by 3.098 million ha and plantation by 386,532 ha. The exported country label footprint is approximately 231.405 million ha. Province and biome totals are respectively about 39,860 and 2,516 ha smaller; their net differences are not geometry symmetric differences or proof of a nested missing-area partition. The country series, not a forced subunit sum, supplies the reported result. The ATBD documents class confusion, temporal irregularities and ecotone integration using stable map labels. No compatible reference-based change interval is established, and the agricultural increase is not attributed to forest conversion from marginal totals alone.

For Brazil, Indonesia and Congo, respectively, symmetric endpoint net-area errors would each have to be below approximately 5.431, 0.351 and 1.652 million ha to guarantee the reported negative forest-group difference under that error model. These are requirements on unknown errors, not measured uncertainty budgets. Cover and biomass remain distinct.

\subsection{Unique mapped endpoint cohorts, not repeated annual events}\label{sec:flowresults}
Brazilian formations have 22.711 million outflow hectare-events and 11.848 million inflow hectare-events during the eight annual transitions; Indonesia has 4.424 and 3.722 million, respectively. Event totals exceed net change and need not equal unique converted area.

The geographically restricted marked-flow networks imply Brazil formations-to-nonforest endpoint cohorts between 11.204 and 22.711 million ha; the corresponding cropland range is 0--7.014 million ha. Adding savannas changes those ranges to 18.374--35.847 and 0.009--8.662 million ha, respectively. Indonesian forest-to-nonforest is bounded by 1.312--4.424 million ha, forest-to-agriculture by 0.127--3.764 million ha, and forest-to-oil-palm by \textbf{0.112--2.326 million ha}. Each range includes its outward numerical allowance and conditions on the annual mapped capacities. Geographic restrictions tighten or preserve all raw bounds; all 634 saved primal/dual certificates satisfy the recorded residual criteria.

The positive Indonesian oil-palm lower bound survives every aggregate-compatible intervening class path, but not arbitrary reference error. Its mapped minimum is 111,612 ha: retaining a positive reference-defined lower bound requires combined initial-forest and final-oil-palm class-set symmetric-difference error below that amount. No measured error budget meeting that condition is claimed. A small net-area error or pooled overall map accuracy does not supply it. Nor does this bound identify a biomass loss, direct clearing event, conversion association or causal effect.

\subsection{Inventory evidence separates stock, period and uncertainty}\label{sec:inventoryresults}
Finnish NFI12 and NFI13 growing-stock volumes are 2,475.192 and 2,551.964 million m$^3$, with reported standard errors 13.001 and 13.534 million m$^3$. Their difference is $+76.772$ million m$^3$ ($+3.10\%$). Equation~\ref{eq:invbonf}, with source rounding, gives an approximate 95\% sampling band of $[17.293,136.251]$ million m$^3$ without cross-cycle independence. The second-moment sensitivity is $[-41.896,195.440]$ million m$^3$, unresolved. An unbounded differential systematic bias is not covered by either result. Total living-tree biomass is 1,744 versus 1,786 Mt including roots and stumps; the sum of stems plus branches/foliage is 1,350 versus 1,384 Mt. Neither receives the volume standard error or is promoted to the original AGB target.

French production-forest living volume is $2,716\pm45$ million m$^3$ in 2013--2017 and $2,916\pm47$ million m$^3$ in 2021--2025, where the published margins are 95\% endpoint halfwidths. The difference is $+200$ million m$^3$ ($+7.36\%$). A dependence-free band with at least 90\% sampling coverage, including full-unit rounding allowances, is $[104,296]$ million m$^3$. Under the additional marginal-normal interpretation, the approximate 95\% band is $[90.502,309.498]$ million m$^3$; four-endpoint simultaneous adjustment gives $[78.210,321.790]$ million m$^3$ for this longest contrast. These remain conditional on comparable measurement targets and methods, with differential bias $B$ widening either side. The shorter 2018--2022 and 2019--2023 comparisons against 2021--2025 are also retained: their approximate 95\% pairwise bands include zero. The latest overlapping windows must not be used as independent replicates, and none is an exact endpoint AGB comparison.

The Canadian inventory-domain annualized AGB change is $+17.639$ Mt/year with a provider-labelled halfwidth of 38.918 Mt/year, giving the unrounded reported range $[-21.279,56.557]$ Mt/year. The net direction is unresolved by that range. Component point changes are $-16.158$ (coniferous), $-16.115$ (mixedwood), $+43.065$ (broadleaf) and $+6.847$ (non-treed) Mt/year; they concern changing class-defined stocks, not tracked biological cohorts or a decomposition of causal growth. All component intervals are retained in the data. The separate total relative standard error is labelled 1.16\%; applying the usual definition would imply a standard error of 0.205 Mt/year, irreconcilable with the 38.918 halfwidth under ordinary normal or large-sample $t$ conventions. The stock table does approximately reconcile, so the change export is not silently treated as the same statistical representation. We do not resolve the discrepancy by selecting the smaller uncertainty or assigning an unsupported confidence level. Projection, artificial-change filtering, geographic omissions and timing further prevent a primary national AGB inference.
